\setcounter{framenumber}{0}
%27

\begin{frame}
  \titlepage
\end{frame}

\begin{frame}{Learning Goals}
By the end of this lecture, students should be able to:
\begin{itemize}
    \item compute distributions of transformed random variables;
    \item apply the change-of-variables formula;
    \item define convolutions;
    \item compute distributions of sums of independent random variables;
\end{itemize}
\end{frame}

\lecturepart{Transformations of Random Variables}

\begin{frame}{Why Transform Random Variables?}
Transformations appear everywhere in probability and statistics.

\vspace{0.25cm}

Examples:
\begin{itemize}
    \item unit conversion;
    \item standardization;
    \item logarithms;
    \item sums and averages;
    \item squares and absolute values;
    \item maxima and minima.
\end{itemize}

\vspace{0.25cm}

\begin{keybox}{Main question}
If
\[
Y=g(X),
\]
what is the distribution of $Y$?
\end{keybox}
\end{frame}






\begin{frame}{Transformations in the Discrete Case}

If $X$ is discrete and
\[
Y=g(X),
\]
then, for every possible value $y$ of $Y$,
\[
{
\Prob(Y=y)
=
\Prob\!\left(X\in g^{-1}(\{y\})\right)
=
\sum_{x:\,g(x)=y}
\Prob(X=x).
}
\]

In words: find all values of $x$ that are mapped to $y$ by $g$, and add their probabilities.

\vspace{0.25cm}

\textbf{Example.}
Suppose
\[
\Prob(X=-2)=0.1,\quad
\Prob(X=-1)=0.2,\quad
\Prob(X=0)=0.3,\quad
\Prob(X=1)=0.25,\quad
\Prob(X=2)=0.15,
\]
and let
\[
Y=X^2.
\]

Then $Y\in\{0,1,4\}$, and
\[
\begin{aligned}
\Prob(Y=0)
&=\Prob(X=0)=0.30,\\
\Prob(Y=1)
&=\Prob(X\in\{-1,1\})
=\Prob(X=-1)+\Prob(X=1)=0.45,\\
\Prob(Y=4)
&=\Prob(X\in\{-2,2\})
=\Prob(X=-2)+\Prob(X=2)=0.25.
\end{aligned}
\]

\end{frame}
















%=========================================================
\begin{frame}{Transformations of Continuous Random Variables}

\small

\begin{keybox}{Change of Variables Theorem (Univariate Case)}
Let $X$ be a continuous random variable with PDF $f_X$ and support $S_X$,
and let
\[
Y=g(X),
\]
where $g$ is differentiable and strictly monotone on $S_X$.

Then $g$ has an inverse, and the support of $Y$ is
\[
S_Y=g(S_X).
\]

For $y\in S_Y$,
\[
{
f_Y(y)
=
f_X\!\left(g^{-1}(y)\right)
\left|
\frac{d}{dy}g^{-1}(y)
\right|
}
\]
and $f_Y(y)=0$ for $y\notin S_Y$.
\end{keybox}
 

\textbf{How to use it:}
solve for $x$ in terms of $y$, find $\left|dx/dy\right|$, substitute into
$f_X(x)$, and determine the support of $Y$.

\textbf{Note:} $x=g^{-1}(y)$, so by $\frac{dx}{dy}$, we mean $\frac{d}{dy}g^{-1}(y)$. 


\begin{keybox}{Important}
The theorem requires $g$ to be one-to-one on the support of $X$. When it is not, we can break the support into monotone cases, or follow the same steps of the proof and adjust the step for isolating $x$ appropriately (examples below).
 
\end{keybox}

\end{frame}


%=========================================================
\begin{frame}{Example 1: A Monotone Transformation}

\small

Let
\[
X\sim N(0,1),
\qquad
Y=e^X.
\]

Since
\[
g(x)=e^x
\]
is strictly increasing on $\mathbb R$, the change-of-variables theorem applies.

\textbf{Step 1: Find the support of $Y$.}

Since $e^x>0$ for every $x$,
\[
S_Y=(0,\infty).
\]

\textbf{Step 2: Solve $y=g(x)$ for $x$.}
\[
y=e^x
\quad\Longrightarrow\quad
x=\log y.
\]

Thus
\[
g^{-1}(y)=\log y,
\qquad
\left|\frac{dx}{dy}\right|=\frac{1}{y}.
\]

\textbf{Step 3: Apply the formula.}

Since
\[
f_X(x)=\frac{1}{\sqrt{2\pi}}e^{-x^2/2},
\]
we obtain
\[
{
f_Y(y)
=f_X\!\left(g^{-1}(y)\right)
\left|
\frac{d}{dy}g^{-1}(y)
\right|=
\frac{1}{y\sqrt{2\pi}}
\exp\left\{-\frac{(\log y)^2}{2}\right\},
\qquad y>0.
}
\]

For $y\le0$, $f_Y(y)=0$.

This is the \textbf{log-normal distribution}.

\end{frame}

%=========================================================
\begin{frame}{Example 2: When the Transformation Is Not One-to-One}

\small

Let
\[
X\sim N(0,1),
\qquad
Y=X^2.
\]

Here
\[
g(x)=x^2
\]
is \textbf{not one-to-one} on $\mathbb R$:
\[
g(-x)=g(x).
\]
Therefore, we cannot define a single inverse $g^{-1}$ on all of $\mathbb R$.
Instead, use the \textbf{CDF method}.

Since $Y=X^2\ge0$, the support is
\[
S_Y=[0,\infty).
\]

For $y>0$,
\begin{equation}
F_Y(y)
=\Prob(Y\le y)
=\Prob(X^2\le y)
=\Prob(-\sqrt y\le X\le\sqrt y)
=\Phi(\sqrt y)-\Phi(-\sqrt y)
=2\Phi(\sqrt y)-1.
\end{equation}

Differentiate:
\begin{equation}
f_Y(y)
=\frac{d}{dy}\left[2\Phi(\sqrt y)-1\right]
=2\varphi(\sqrt y)\frac{1}{2\sqrt y}
=\frac{\varphi(\sqrt y)}{\sqrt y},
\qquad y>0.
\end{equation}

Therefore, using
\[
\varphi(x)=\frac{1}{\sqrt{2\pi}}e^{-x^2/2},
\]
the PDF of $Y$ is
\[
{
f_Y(y)
=
\begin{cases}
\displaystyle
\frac{1}{\sqrt{2\pi y}}e^{-y/2},
& y>0,\\[6pt]
0,
& y\le 0.
\end{cases}
}
\]

\end{frame}




%=========================================================
\begin{frame}{Transformations of Continuous Random Vectors}

\small

\begin{keybox}{Change of Variables Theorem (Multivariate Case)}
Let $\mathbf X=(X_1,\ldots,X_n)^T$ be a continuous random vector with joint PDF $f_{\mathbf X}$ and support
$S_{\mathbf X}\subseteq\mathbb R^n$. Define $\mathbf Y=g(\mathbf X),
$ where $g:S_{\mathbf X}\to S_{\mathbf Y}$ is one-to-one and differentiable,
with differentiable inverse $\mathbf x=g^{-1}(\mathbf y).
$. Then the support of $\mathbf Y$ is $S_{\mathbf Y}=g(S_{\mathbf X}),
$ and
\[
{
f_{\mathbf Y}(\mathbf y)
=
f_{\mathbf X}\!\left(g^{-1}(\mathbf y)\right)
\left|
\det\!\left(
\frac{\partial \mathbf x}{\partial \mathbf y}
\right)
\right|,
\qquad \mathbf y\in S_{\mathbf Y}.
}
\]

Here
\[
\frac{\partial \mathbf x}{\partial \mathbf y}
=
\begin{pmatrix}
\dfrac{\partial x_1}{\partial y_1} & \cdots &
\dfrac{\partial x_1}{\partial y_n}\\
\vdots & \ddots & \vdots\\
\dfrac{\partial x_n}{\partial y_1} & \cdots &
\dfrac{\partial x_n}{\partial y_n}
\end{pmatrix}
\]
is the \textbf{Jacobian matrix}. For $\mathbf y\notin S_{\mathbf Y}$,
$f_{\mathbf Y}(\mathbf y)=0$.
\end{keybox}

\end{frame}


%=========================================================
\begin{frame}{Example: A Transformation of Two Variables}

\small

Let $X_1,X_2$ be independent $\operatorname{Exp}(1)$ random variables, so
\[
f_{X_1,X_2}(x_1,x_2)
=
e^{-(x_1+x_2)},
\qquad x_1>0,\ x_2>0.
\]

Define
\[
Y_1=X_1+X_2,
\qquad
Y_2=\frac{X_1}{X_1+X_2}.
\]

\textbf{Question: What is $f_{Y_1,Y_2(y_1,y_2)}$}?

Since
\[
y_2=\frac{x_1}{x_1+x_2}=\frac{x_1}{y_1},
\]
we obtain
\[
{
x_1=y_1y_2,
\qquad
x_2=y_1(1-y_2).
}
\]

We now find the support:
 
Since $x_1>0$ and $x_2>0$, we have $y_1y_2>0,
\qquad
y_1(1-y_2)>0,$ which gives
\[
{
y_1>0,
\qquad
0<y_2<1.
}
\]
We now find the Jacobian: 

\[
\frac{\partial(x_1,x_2)}{\partial(y_1,y_2)}
=
\begin{pmatrix}
y_2 & y_1\\
1-y_2 & -y_1
\end{pmatrix},
\]
so
\[
\left|
\det\left(
\frac{\partial(x_1,x_2)}{\partial(y_1,y_2)}
\right)
\right|
=
|-y_1|
=
y_1.
\]

\end{frame}


%=========================================================
\begin{frame}{Example: Finding the Joint PDF}

\small

By the change of variables theorem,
\[
f_{Y_1,Y_2}(y_1,y_2)
=
f_{X_1,X_2}
\left(
y_1y_2,\,
y_1(1-y_2)
\right)
\left|
\frac{\partial(x_1,x_2)}
{\partial(y_1,y_2)}
\right|.
\]

Since
\[
x_1+x_2
=
y_1y_2+y_1(1-y_2)
=
y_1,
\]
we have
\[
f_{X_1,X_2}
\left(
y_1y_2,\,
y_1(1-y_2)
\right)
=
e^{-y_1}.
\]

Therefore,
\[
{
f_{Y_1,Y_2}(y_1,y_2)
=
\begin{cases}
y_1e^{-y_1},
& y_1>0,\quad 0<y_2<1,\\[4pt]
0,
& \text{otherwise}.
\end{cases}
}
\]

Notice that
\[
f_{Y_1,Y_2}(y_1,y_2)
=
\underbrace{y_1e^{-y_1}}_{f_{Y_1}(y_1)}
\underbrace{1}_{f_{Y_2}(y_2)}.
\]

Hence,
\[
{
Y_1\sim\operatorname{Gamma}(2,1),
\qquad
Y_2\sim\operatorname{Unif}(0,1),
\qquad
Y_1\perp Y_2.
}
\]

\end{frame}







\lecturepart{Convolutions}

\begin{frame}{What Is a Convolution?}

A convolution is the distribution of a sum of independent random variables.

If
\[
T=X+Y,
\]
and $X,Y$ are independent, we want the distribution of $T$.

\vspace{0.25cm}

\begin{keybox}{Why sums matter}
Sums appear everywhere:
\begin{itemize}
    \item totals;
    \item averages;
    \item waiting times;
    \item counting processes.
\end{itemize}
\end{keybox}

\end{frame}


%=========================================================
\begin{frame}{Discrete Convolution Formula}

\begin{keybox}{Theorem}
If $X$ and $Y$ are independent discrete random variables and
\[
T=X+Y,
\]
then
\[
{
\Prob(T=t)
=
\sum_x
\Prob(X=x)\Prob(Y=t-x).
}
\]
\end{keybox}

\textbf{Proof.}

For a fixed value $t$, the event
\[
\{X+Y=t\}
\]
can be written as the disjoint union
\[
\{X+Y=t\}
=
\bigcup_x
\{X=x,\;Y=t-x\}.
\]

Therefore,
\[
\Prob(T=t)
=
\sum_x
\Prob(X=x,\;Y=t-x).
\]

Since $X$ and $Y$ are independent,
\[
\Prob(X=x,\;Y=t-x)
=
\Prob(X=x)\Prob(Y=t-x).
\]

Hence,
\[
{
\Prob(T=t)
=
\sum_x
\Prob(X=x)\Prob(Y=t-x).
}
\]

\end{frame}


%=========================================================
\begin{frame}{Continuous Convolution Formula}

\begin{keybox}{Theorem}
If $X$ and $Y$ are independent continuous random variables and
\[
T=X+Y,
\]
then
\[
{
f_T(t)
=
\int_{-\infty}^{\infty}
f_X(x)f_Y(t-x)\,dx.
}
\]
\end{keybox}

\end{frame}








%=========================================================
\begin{frame}{Proof of the Continuous Convolution Formula}

\small

Since
\[
T=X+Y,
\]
the CDF of $T$ is
\[
F_T(t)
=
\Prob(T\le t)
=
\Prob(X+Y\le t).
\]

Using the joint density,
\[
F_T(t)
=
\iint_{\{(x,y):\,x+y\le t\}}
f_{X,Y}(x,y)\,dx\,dy.
\]

For each fixed $x$, the condition
\[
x+y\le t
\]
is equivalent to
\[
y\le t-x.
\]

Therefore,
\[
F_T(t)
=
\int_{-\infty}^{\infty}
\int_{-\infty}^{\,t-x}
f_{X,Y}(x,y)\,dy\,dx.
\]

Since $X$ and $Y$ are independent,
\[
f_{X,Y}(x,y)=f_X(x)f_Y(y),
\]
so
\[
F_T(t)
=
\int_{-\infty}^{\infty}
f_X(x)
\left[
\int_{-\infty}^{\,t-x}
f_Y(y)\,dy
\right]dx.
\]

Differentiating with respect to $t$,
\[
{
f_T(t)
=
\int_{-\infty}^{\infty}
f_X(x)f_Y(t-x)\,dx.
}
\]

\end{frame}




%=========================================================
\begin{frame}{Proof of Convolution Using the Jacobian Formula}

\small

Let
\[
T=X+Y,
\qquad
U=X.
\]

Then the transformation is
\[
t=x+y,
\qquad
u=x.
\]

Solving for the original variables gives
\[
{
x=u,
\qquad
y=t-u.
}
\]

The Jacobian of the inverse transformation is
\[
\frac{\partial(x,y)}{\partial(t,u)}
=
\begin{pmatrix}
\dfrac{\partial x}{\partial t} &
\dfrac{\partial x}{\partial u}\\[6pt]
\dfrac{\partial y}{\partial t} &
\dfrac{\partial y}{\partial u}
\end{pmatrix}
=
\begin{pmatrix}
0 & 1\\
1 & -1
\end{pmatrix},
\]
and therefore
\[
\left|
\det\left(
\frac{\partial(x,y)}{\partial(t,u)}
\right)
\right|
=1.
\]

By the change-of-variables formula,
\[
f_{T,U}(t,u)
=
f_{X,Y}(u,t-u).
\]

Since $X$ and $Y$ are independent,
\[
f_{X,Y}(u,t-u)
=
f_X(u)f_Y(t-u).
\]

Hence,
\[
{
f_{T,U}(t,u)
=
f_X(u)f_Y(t-u).
}
\]

\end{frame}


%=========================================================
\begin{frame}{Proof of Convolution Using the Jacobian Formula (continued)}

\small

We have found the joint PDF
\[
f_{T,U}(t,u)
=
f_X(u)f_Y(t-u).
\]

To obtain the marginal PDF of
\[
T=X+Y,
\]
integrate out $U$:
\[
\begin{aligned}
f_T(t)
&=
\int_{-\infty}^{\infty}
f_{T,U}(t,u)\,du\\
&=
\int_{-\infty}^{\infty}
f_X(u)f_Y(t-u)\,du.
\end{aligned}
\]

Therefore,
\[
{
f_T(t)
=
\int_{-\infty}^{\infty}
f_X(u)f_Y(t-u)\,du.
}
\]

Renaming the dummy variable $u$ as $x$ gives the previous convolution formula:
\[
{
f_T(t)
=
\int_{-\infty}^{\infty}
f_X(x)f_Y(t-x)\,dx.
}
\]

\end{frame}



%=========================================================
\begin{frame}{Example 1: Sum of Two Exponential Random Variables}

\small

Let
\[
X,Y\overset{\text{i.i.d.}}{\sim}\operatorname{Exp}(\lambda),
\qquad
T=X+Y.
\]

Recall that
\[
f_X(x)=f_Y(x)
=
\begin{cases}
\lambda e^{-\lambda x}, & x>0,\\
0, & \text{otherwise}.
\end{cases}
\]

Since $X,Y>0$, we immediately know
\[
T=X+Y>0.
\]

For independent continuous random variables, the convolution formula is
\[
f_T(t)
=
\int_{-\infty}^{\infty}
f_X(x)f_Y(t-x)\,dx.
\]

For the integrand to be nonzero, we need \emph{both}
\[
x>0
\qquad\text{and}\qquad
t-x>0.
\]

Thus, $0<x<t,$ which is possible only when $t>0$.

Therefore, for $t>0$,
 
\begin{equation}
f_T(t)
=
\int_0^t
f_X(x)f_Y(t-x)\,dx =
\int_0^t
\lambda e^{-\lambda x}
\lambda e^{-\lambda(t-x)}
\,dx =
\lambda^2e^{-\lambda t}\int_0^t dx =
\lambda^2te^{-\lambda t}.
\end{equation}
 

Hence,
\[
f_T(t)
=
\begin{cases}
\lambda^2te^{-\lambda t}, & t>0,\\
0, & t\le0.
\end{cases}
\]
 
Thus, $T=X+Y\sim\operatorname{Gamma}(2,\lambda).$
 

\end{frame}


%=========================================================
\begin{frame}{Example 2: Sum of Two Uniform Random Variables}

\small

Let
\[
X,Y\overset{\text{i.i.d.}}{\sim}\operatorname{Unif}(0,1),
\qquad
T=X+Y.
\]

Recall that
\[
f_X(x)=f_Y(x)
=
\begin{cases}
1, & 0<x<1,\\
0, & \text{otherwise}.
\end{cases}
\]

Since
\[
0<X<1,
\qquad
0<Y<1,
\]
we know
\[
0<T<2.
\]

The convolution formula gives
\[
f_T(t)
=
\int_{-\infty}^{\infty}
f_X(x)f_Y(t-x)\,dx.
\]

The integrand is nonzero exactly when
\[
0<x<1
\qquad\text{and}\qquad
0<t-x<1.
\]

The second inequality gives
\[
t-1<x<t.
\]

Hence $x$ must satisfy both
\[
0<x<1
\qquad\text{and}\qquad
t-1<x<t.
\]

Therefore,
\[
\max(0,t-1)<x<\min(1,t).
\]

So the integration bounds depend on $t$.

\end{frame}


%=========================================================
\begin{frame}{Example 2: Determining the PDF}

\small

We found that
\[
f_T(t)
=
\int_{\max(0,t-1)}^{\min(1,t)}1\,dx,
\qquad 0<t<2.
\]

Now consider the two possible regions.

\textbf{Case 1: $0<t\le1$.}

Then
\[
\max(0,t-1)=0,
\qquad
\min(1,t)=t,
\]
so
\[
f_T(t)
=
\int_0^t 1\,dx
=
t.
\]

\textbf{Case 2: $1<t<2$.}

Then
\[
\max(0,t-1)=t-1,
\qquad
\min(1,t)=1,
\]
so
\[
f_T(t)
=
\int_{t-1}^1 1\,dx
=
2-t.
\]

Therefore,
\[
f_T(t)
=
\begin{cases}
t, & 0<t\le1,\\[2mm]
2-t, & 1<t<2,\\[2mm]
0, & \text{otherwise}.
\end{cases}
\]

This is a \textbf{triangular distribution}: values near $1$ can be obtained
from more pairs $(X,Y)$ than values near $0$ or $2$.

\end{frame}






 